\chapter[Dự đoán giá xe máy cũ]{Dự đoán giá xe máy cũ ở TP.HCM và Hà Nội}

\section{Bài toán và hướng tiếp cận}

Người mua và người bán xe máy cũ thường không biết một chiếc xe ``đáng giá'' bao nhiêu: cùng là Honda Vision, xe đời
2016 và đời 2023 chênh nhau vài lần, xe chạy 5 000 km khác xe chạy 80 000 km, và cùng một chiếc xe có thể được rao giá
khác nhau ở hai thành phố. Bài tập này xây dựng mô hình ước lượng \textbf{giá rao bán} của một chiếc xe máy cũ từ các
thông tin có trong tin đăng (hãng, dòng xe, đời xe, số km, nơi bán, người bán), đồng thời trả lời ba câu hỏi:

\begin{itemize}
  \item Giá xe giảm bao nhiêu theo đời xe và số km?
  \item Cùng một dòng xe, cùng đời, giá ở Hà Nội có khác TP.HCM không, và vì sao?
  \item Dữ liệu tự thu thập có những lỗi gì, và cần xử lý chúng thế nào?
\end{itemize}

Đây là bài toán \textbf{hồi quy có giám sát}. Quy trình gồm các bước: kiểm tra chất lượng dữ liệu và giới hạn bài toán
vào các dòng xe phổ biến; chia train / validation / test; phân tích khám phá trên train; tạo feature, điền giá trị
thiếu, mã hóa và chuẩn hóa trong một \texttt{Pipeline} của scikit-learn \cite{tab-sklearn}; so sánh hai baseline với năm
mô hình hồi quy bằng cross-validation và validation; chấm trên tập test và phân tích lỗi. Notebook chạy được trên Google
Colab với \textit{Runtime > Run all} trong khoảng một phút, không cần GPU.

\section{Bộ dữ liệu}

\begin{itemize}
  \item \textbf{Nguồn:} chuyên mục \textit{Xe máy} của Chợ Tốt \cite{tab-chotot}, trang rao vặt lớn ở Việt Nam. Nhóm tự
        thu thập bằng script \texttt{crawl\_chotot\_xemay.py} qua API công khai mà chính trang web dùng để hiển thị danh
        sách tin, không cần đăng nhập.
  \item \textbf{Cách thu thập:} API trả tối đa 50 tin mỗi trang và 10 000 tin mỗi truy vấn, nên script lấy lần lượt
        từng quận/huyện của TP.HCM và Hà Nội, khoảng một request mỗi giây. Hãng, dòng xe, loại xe, xuất xứ và dung tích
        được API trả về dạng mã số; script tra nhãn chữ của từng mã bằng API chi tiết của một tin có mã đó.
  \item \textbf{Thời điểm:} thu thập ngày 16/11/2026. Mỗi dòng là một tin đang hiển thị lúc đó, đăng từ 17/09 đến
        16/11/2026.
  \item \textbf{Ẩn danh:} file CSV đã bỏ tên người bán, ID tài khoản, ảnh đại diện và phần mô tả tự do; số điện thoại
        trong tiêu đề được thay bằng \texttt{[SĐT]}. Chỉ giữ các con số tổng hợp về người bán (số tin đang rao, số tin
        đã bán, điểm đánh giá). Dữ liệu chỉ dùng cho mục đích học tập.
  \item \textbf{Lưu ý về nhãn:} giá trong dữ liệu là giá rao do người bán tự đặt, không phải giá giao dịch thật.
\end{itemize}

Bộ dữ liệu có 19 718 tin (15 890 ở TP.HCM, 3 828 ở Hà Nội) và 28 cột, đáp ứng mọi yêu cầu của đề với dữ liệu dạng bảng
(Bảng~\ref{tab:tab-req}). Bảng~\ref{tab:tab-cols} mô tả các nhóm cột.

\begin{table}[H]\centering
\caption{Bộ dữ liệu so với yêu cầu của đề (mục 4.2).}\label{tab:tab-req}
\begin{tabular}{ll}
\toprule
Yêu cầu & Bộ dữ liệu \\
\midrule
Ít nhất 2 000 dòng & 19 718 dòng \\
Ít nhất 10 cột & 28 cột \\
Có giá trị thiếu & 13 cột có ô thiếu \\
Có cột phân loại & 8 cột, ví dụ \texttt{model} có 403 mức \\
Có cột số & 7 cột, ví dụ \texttt{mileage\_km}, \texttt{reg\_year} \\
Có độ khó & giá lệch phải (skew 71.5), 7.1\% tin ngoài khoảng IQR \\
\bottomrule
\end{tabular}
\end{table}

\begin{table}[H]\centering
\caption{Các nhóm cột của bộ dữ liệu.}\label{tab:tab-cols}
\begin{tabular}{>{\raggedright\arraybackslash}p{2.6cm}>{\raggedright\arraybackslash}p{6.2cm}>{\raggedright\arraybackslash}p{5.4cm}}
\toprule
Nhóm & Cột & Ý nghĩa \\
\midrule
Mã và metadata & \texttt{list\_id}, \texttt{subject}, \texttt{posted\_at}, \texttt{orig\_posted\_at}, \texttt{crawled\_at} & mã tin, tiêu đề, ngày đăng, ngày thu thập \\
Nơi bán & \texttt{city}, \texttt{district}, \texttt{ward}, \texttt{latitude}, \texttt{longitude} & thành phố, quận/huyện, phường/xã, tọa độ \\
Chiếc xe & \texttt{brand}, \texttt{model}, \texttt{bike\_type}, \texttt{origin}, \texttt{capacity}, \texttt{reg\_year}, \texttt{mileage\_km}, \texttt{condition} & hãng, dòng xe, loại xe, xuất xứ, dung tích, năm đăng ký, số km, tình trạng \\
Người bán & \texttt{seller\_is\_shop}, \texttt{seller\_live\_ads}, \texttt{seller\_sold\_ads}, \texttt{seller\_avg\_rating}, \texttt{seller\_total\_rating} & cửa hàng hay cá nhân, số tin đang rao, đã bán, điểm và số lượt đánh giá \\
Tin đăng & \texttt{number\_of\_images}, \texttt{has\_video}, \texttt{inspected}, \texttt{is\_price\_not\_valid} & số ảnh, video, kiểm định, cờ giá không hợp lệ \\
\textbf{Target} & \texttt{price\_vnd} & giá rao bán (VND) \\
\bottomrule
\end{tabular}
\end{table}

\section{Kiểm tra chất lượng dữ liệu}

Dữ liệu rao vặt có nhiều lỗi do người đăng tự nhập. Các luật sửa lỗi ở mục này là luật cố định theo từng dòng, không
học con số nào từ dữ liệu, nên được áp dụng trên toàn bảng trước khi chia tập mà không gây rò rỉ. Bảng~\ref{tab:tab-quality}
tổng hợp kết quả.

\begin{table}[H]\centering
\caption{Các vấn đề chất lượng dữ liệu và cách xử lý.}\label{tab:tab-quality}
\begin{tabular}{>{\raggedright\arraybackslash}p{4.4cm}>{\raggedright\arraybackslash}p{5.2cm}>{\raggedright\arraybackslash}p{4.6cm}}
\toprule
Vấn đề & Kết quả kiểm tra & Cách xử lý \\
\midrule
Giá trị thiếu & \texttt{origin} thiếu 56.1\%, \texttt{capacity} 43.4\% (gồm 132 giá trị ẩn \textit{Không biết rõ}, \textit{mã 0}), \texttt{mileage\_km} 1.2\% & điền trong pipeline (mục \ref{sec:tab-prep}) \\
Cột hằng & \texttt{is\_price\_not\_valid}, \texttt{has\_video} chỉ có một giá trị; \texttt{inspected} cũng vậy, trừ 1 ô thiếu & bỏ cột \\
Tọa độ ở mức phường & 1 243 cặp tọa độ cho 19 718 tin, gần bằng số phường (428) & bỏ tọa độ và phường, giữ quận \\
Giá rác & 1 tin thiếu giá, 42 tin dưới 2 triệu, 2 tin trên 2 tỷ (một tin thu mua 8.9 tỷ, một chiếc Cub 6.5 tỷ) & bỏ tin \\
Tin cần mua / thu mua & 32 tin, ví dụ ``MUA XE CŨ TẬN NHÀ GIÁ CAO'' & bỏ tin \\
Năm đăng ký không thể & 3 tin trước 1980 hoặc sau năm đăng tin & bỏ tin \\
Số km ``ảo'' & 2 093 tin xe đã dùng (10.8\%) khai 0, 1, 12 345, 99 999 km...; 168 tin trên 500 000 km & coi là thiếu, giữ tin \\
Tin đăng lại & 518 tin (2.6\%) trùng thành phố, quận, hãng, dòng, năm, số km và giá & giữ bản đầu tiên \\
\bottomrule
\end{tabular}
\end{table}

\reportfig{fig02_fake_km.png}{Mười lăm số km được khai nhiều nhất trong tin xe đã dùng. Màu cam là giá trị không thể
có: 12 345 km xuất hiện 358 lần, nhiều hơn cả 12 000 km (352 lần). Các giá trị này được coi là thiếu chứ không phải số
km thật.}{fig:tab-fakekm}

Ngưỡng trên 2 tỷ chứ không phải 1 tỷ vì hai tin đắt tiếp theo là Harley-Davidson CVO khoảng 1.5 tỷ, giá thật của dòng
xe này. Có 54 tin ghi đúng năm 1980, nhiều hơn hẳn các năm 1981--1989, nhiều khả năng là lựa chọn ``1980 trở về trước''
của trang, nên vẫn được giữ.

\subsection{Phạm vi bài toán: các dòng xe phổ biến}

Sau khi làm sạch, dữ liệu có hơn 400 tổ hợp hãng + dòng xe, nhưng phần lớn chỉ có vài tin. Với dòng xe chỉ có 3--5
tin, mô hình không có đủ ví dụ để học, và điểm trên test của dòng đó cũng không đáng tin. Vì vậy bài toán được giới hạn
vào các dòng xe có \textbf{ít nhất 40 tin}; khi chia 70 / 15 / 15, mỗi dòng có khoảng 28 tin để học và 6 tin để chấm.
Nhóm \texttt{Dòng khác} luôn nằm ngoài phạm vi vì là nhóm gộp của nhiều loại xe không rõ tên. Luật chỉ dựa trên số tin
của mỗi dòng, không dựa trên giá của từng tin, và hàm dự đoán từ chối xe ngoài phạm vi.

\reportfig{fig03_scope.png}{Tỉ lệ tin được phủ theo số dòng xe. 65 dòng xe có từ 40 tin chiếm 82.8\% số tin; 343 dòng
còn lại chỉ chiếm 11.6\% (trung bình 6.5 tin mỗi dòng), nhóm \texttt{Dòng khác} chiếm 5.6\%.}{fig:tab-scope}

Bảng~\ref{tab:tab-funnel} cho biết số tin sau từng bước. Dữ liệu còn lại được chia 70 / 15 / 15 thành train (11 088
tin), validation (2 376) và test (2 377), phân tầng theo thành phố và 10 khoảng giá để ba tập có cùng tỉ lệ Hà Nội
(19\%) và cùng giá median (19.0 triệu). Tập được chia \textbf{trước} phân tích khám phá, để mọi quyết định sau đó chỉ dựa
trên train.

\begin{table}[H]\centering
\caption{Số tin sau từng bước làm sạch.}\label{tab:tab-funnel}
\begin{tabular}{lrr}
\toprule
Bước & Số tin & Bị bỏ \\
\midrule
Dữ liệu thô & 19 718 & \\
Bỏ giá rác, tin cần mua, năm không thể & 19 644 & 74 \\
Bỏ tin đăng lại & 19 127 & 517 \\
Giữ dòng xe phổ biến & 15 841 & 3 286 \\
\bottomrule
\end{tabular}
\end{table}

\section{Phân tích khám phá dữ liệu}

Phân tích khám phá chỉ dùng tập train.

\reportfig{fig04_target.png}{Phân phối giá rao trên thang gốc và thang log. Giá lệch phải mạnh (skew 3.3) nhưng gần đối
xứng trên thang log (skew 0.03); giá median ở TP.HCM là 20.5 triệu, ở Hà Nội 14.5 triệu. Vì vậy mô hình học trên
$\log_{10}$(giá), để sai 10\% ở xe rẻ và ở xe đắt được tính như nhau.}{fig:tab-target}

\reportfig{fig05_price_by_model.png}{Giá theo 15 dòng xe nhiều tin nhất. Median của SH (68 triệu) gấp 9 lần Sirius (7.6
triệu): dòng xe gần như quyết định mức giá.}{fig:tab-model}

Train có 65 tổ hợp hãng + dòng xe và 50 quận/huyện, bằng hoặc vượt ngưỡng 50 mức mà bài EDA mẫu của môn học gợi ý dùng
\textbf{target encoding} thay cho one-hot. Mục~\ref{sec:tab-enc} kiểm tra lựa chọn này bằng thí nghiệm.

\reportfig{fig06_depreciation.png}{Giá median theo năm đăng ký (từ 2013) và theo số km của 4 dòng xe nhiều tin nhất. Mỗi
năm đời xe cũ hơn, giá giảm khoảng 9\% với Air Blade, 7\% với SH, 5\% với Wave và 7\% với Vision. Xe đi trên 80 000 km
còn 61\%, 64\%, 88\% và 74\% giá xe đi dưới 5 000 km (cùng thứ tự).}{fig:tab-dep}

Giá giảm gần như đều theo năm trên thang log, còn theo số km thì giảm không đều vì số km đi cùng đời xe (xe cũ thường đi
nhiều hơn). Vì vậy bước tiền xử lý thêm feature \textbf{tuổi xe} và \textbf{số km mỗi năm} để tách hai ảnh hưởng. Đường
mất giá chỉ vẽ từ 2013 vì SH đăng ký năm 2010--2012 phần lớn là bản nhập khẩu từ Ý, đắt hơn bản lắp ráp sau đó, làm SH
trông như ``không mất giá''.

Để so sánh giá giữa hai thành phố một cách công bằng, mỗi tin được chia cho giá median của các tin TP.HCM cùng dòng xe
và cùng năm đăng ký (chỉ dùng nhóm có ít nhất 5 tin TP.HCM). Tỉ số bằng 1 nghĩa là rao bằng giá TP.HCM.

\reportfig{fig07_hanoi_vs_hcm.png}{Giá rao so với giá median TP.HCM cùng dòng xe và cùng đời. Người bán cá nhân ở Hà Nội
rao bằng 0.90 lần giá TP.HCM, nhưng cửa hàng chỉ 0.62 lần, thấp nhất ở Quận Nam Từ Liêm (0.50 lần); quận này chiếm 32\%
số tin Hà Nội. Ở TP.HCM, cá nhân và cửa hàng rao như nhau.}{fig:tab-city}

Như vậy phần lớn chênh lệch giữa hai thành phố đến từ một nhóm cửa hàng ở Hà Nội rao giá rất thấp, không phải từ khác
biệt chung của thị trường. Giá rao của nhóm này có thể không phải giá bán thật (ví dụ giá mồi để khách gọi điện), điều
dữ liệu hiện có chưa kiểm chứng được. Mô hình cần biết thành phố, quận và loại người bán, nên giữ cả ba cột.

Với các cột số, \texttt{reg\_year} tương quan với $\log$(giá) mạnh nhất (Spearman +0.44), tiếp theo là
\texttt{number\_of\_images} (+0.32); không cột số nào vượt 0.5, nên phần lớn khác biệt về giá nằm ở dòng xe. Cặp cột số
tương quan mạnh nhất là \texttt{seller\_total\_rating} và \texttt{seller\_live\_ads} (+0.52), chưa tới ngưỡng 0.7 của bài
EDA mẫu, nên giữ mọi cột. Theo quy tắc IQR, 6.9\% giá ở thang gốc là outlier, nhưng đó là các xe phân khối lớn có thật;
trên thang log chỉ còn 0.3\%. Vì vậy không bỏ tin nào theo IQR.

\section{Tiền xử lý dữ liệu}\label{sec:tab-prep}

Mọi thống kê học từ dữ liệu (ngưỡng giá, median để điền, mean và std để chuẩn hóa, giá trung bình theo dòng xe) chỉ được
fit trên train rồi dùng lại cho validation và test. Các bước theo thứ tự:

\begin{enumerate}
  \item \textbf{Bỏ tin giá thấp vô lý:} tin có giá dưới 15\% giá median của dòng xe đó (median tính trên train, chỉ dùng
        cho dòng có ít nhất 20 tin), thường là tiền trả trước hoặc tiền cọc, ví dụ một chiếc SH 2025 rao 10 triệu. Luật
        bỏ 12 tin ở train, 2 ở validation và 3 ở test; còn lại 11 076 / 2 374 / 2 374 tin.
  \item \textbf{Tạo feature} theo từng dòng (không học gì từ dữ liệu):
  \begin{itemize}
    \item \texttt{age} (tuổi xe khi đăng tin), \texttt{log\_km} và \texttt{km\_per\_year} sau khi coi số km ảo là thiếu.
          Nhờ vậy 11.3\% số tin train bị coi là thiếu số km, trong khi cột gốc chỉ thiếu 0.7\%.
    \item \texttt{cc\_est}: dung tích lấy từ tiêu đề (``Vision 110'', ``SH 150i''), nếu không có thì lấy giữa khoảng của
          cột \texttt{capacity}. Tiêu đề cho dung tích ở 45.7\% số tin, nên \texttt{cc\_est} chỉ còn thiếu 30.7\% so
          với 42.3\% của \texttt{capacity}.
    \item \texttt{kw\_*}: tiêu đề có nhắc tới ABS, smartkey, ``zin'', ``lướt'', biển số đẹp, hàng nhập, xe độ hay trả
          góp không.
    \item \texttt{brand\_model} và \texttt{model\_cc} (hãng + dòng xe + dung tích), vì cùng tên ``SH'' có bản 125, 150
          và 350 với giá rất khác nhau.
  \end{itemize}
  \item \textbf{Điền giá trị thiếu, mã hóa, chuẩn hóa} trong một \texttt{ColumnTransformer} đặt đầu \texttt{Pipeline}
        của mỗi mô hình (Bảng~\ref{tab:tab-pipe}). Sau bước này, 27 feature gốc thành 167 cột.
\end{enumerate}

\begin{table}[H]\centering
\caption{Xử lý từng nhóm cột trong \texttt{ColumnTransformer}.}\label{tab:tab-pipe}
\begin{tabular}{>{\raggedright\arraybackslash}p{2.6cm}>{\raggedright\arraybackslash}p{4.6cm}>{\raggedright\arraybackslash}p{2.9cm}>{\raggedright\arraybackslash}p{4.1cm}}
\toprule
Nhóm & Cột & Điền thiếu & Mã hóa / chuẩn hóa \\
\midrule
Số & \texttt{age}, \texttt{log\_km}, \texttt{km\_per\_year}, \texttt{cc\_est}, số ảnh, các cột người bán & median của train, kèm cột cờ ``bị thiếu'' & \texttt{StandardScaler} \\
Phân loại ít mức & \texttt{city}, \texttt{bike\_type}, \texttt{origin}, \texttt{capacity}, \texttt{condition} & nhóm \texttt{không rõ} & one-hot, mức dưới 20 tin gộp thành ``hiếm'' \\
Phân loại nhiều mức & \texttt{brand}, \texttt{brand\_model}, \texttt{model\_cc}, \texttt{district} & coi thiếu là một mức & \texttt{TargetEncoder} rồi \texttt{StandardScaler}, thêm one-hot cho \texttt{brand}, \texttt{brand\_model}, \texttt{district} (mức từ 5 tin) \\
Cờ 0/1 & \texttt{seller\_is\_shop}, \texttt{kw\_*} & không thiếu & giữ nguyên \\
\bottomrule
\end{tabular}
\end{table}

\textbf{Target encoding} \cite{tab-te} thay mỗi mức (ví dụ \texttt{Honda | SH}) bằng giá trung bình (thang log) của các
xe thuộc mức đó trong train, có làm mượt: mức có ít tin bị kéo về gần giá trung bình chung. Vì dùng chính target, lớp
\texttt{TargetEncoder} của scikit-learn dùng cross-fitting khi \texttt{fit\_transform} trên train: giá trị của mỗi dòng
được tính từ 4/5 train còn lại, nên mô hình không ``thấy trước'' đáp án. Ví dụ, SH (998 tin) được mã hóa thành 65.1 triệu,
gần median thật 68.0 triệu; Vision (831 tin) thành 18.8 triệu so với median 19.8 triệu.

\reportfig{fig08_scaling.png}{Ba cột số trước và sau \texttt{StandardScaler} (fit trên train). Trước chuẩn hóa,
\texttt{age} có mean 9.3, std 6.5 còn \texttt{seller\_live\_ads} có mean 40.0, std 81.3; sau chuẩn hóa cả ba có mean 0 và
std 1. Hình dạng phân phối không đổi (\texttt{seller\_live\_ads} vẫn lệch phải), nên chuẩn hóa không thay cho việc lấy
log.}{fig:tab-scale}

\section{Xây dựng mô hình}

\textbf{Thước đo chính} là $R^2$ trên $\log_{10}$(giá), đo sai số theo tỉ lệ chứ không theo số tiền. Kèm theo đó là các
thước đo dễ hiểu: sai số tuyệt đối theo \% của giá thật, lấy trung vị trên các tin (\textit{sai số \% median}); tỉ lệ
tin đoán lệch không quá 10\%; và MAE tính bằng triệu đồng. Mỗi mô hình được chấm bằng cross-validation 5 fold trên train
(\texttt{KFold}, xáo trộn, seed 261), rồi huấn luyện lại trên toàn train và chấm trên validation. Mô hình có $R^2$
validation cao nhất được chọn; hai mô hình chỉ được coi là khác nhau khi hiệu $R^2$ lớn hơn độ lệch giữa các fold.

Hai baseline cho biết thế nào là ``đủ tốt'': \texttt{DummyRegressor} luôn đoán giá trung bình của train, và baseline thứ
hai đoán giá median của dòng xe trên train, giống cách một người mua thường ước lượng (``SH cũ khoảng 60--70 triệu''). Các mô hình
trong slide của môn học được so sánh trên cùng bộ tiền xử lý: hồi quy tuyến tính, kNN, cây quyết định, random forest
\cite{tab-rf}. Gradient boosting (\texttt{HistGradientBoostingRegressor}, phần đọc thêm của slide cây và ensemble) được
thêm như một mô hình mở rộng, thử ba cấu hình và giữ cấu hình tốt nhất theo validation (learning rate 0.03, 63 lá,
1 500 vòng).

\begin{table}[H]\centering
\caption{Kết quả cross-validation trên train và trên validation.}\label{tab:tab-val}
\begin{tabular}{lrrr}
\toprule
Mô hình & $R^2$ CV (± độ lệch) & $R^2$ validation & Sai số \% median \\
\midrule
Baseline: giá trung bình & 0.000 ± 0.000 & 0.000 & 49.6\% \\
Baseline: median theo dòng xe & 0.566 ± 0.016 & 0.579 & 30.0\% \\
Hồi quy tuyến tính & 0.753 ± 0.012 & 0.749 & 21.3\% \\
kNN ($k=10$) & 0.771 ± 0.011 & 0.768 & 19.7\% \\
Cây quyết định (sâu 12) & 0.783 ± 0.014 & 0.788 & 15.4\% \\
\textbf{Random forest} & \textbf{0.859 ± 0.011} & \textbf{0.858} & \textbf{12.1\%} \\
Gradient boosting & 0.865 ± 0.008 & 0.856 & 13.2\% \\
\bottomrule
\end{tabular}
\end{table}

Random forest và gradient boosting vượt xa các mô hình còn lại, nhưng thứ tự giữa hai mô hình không ổn định: trên
cross-validation gradient boosting cao hơn 0.006, trên validation random forest cao hơn 0.002. Cả hai hiệu đều nhỏ hơn độ
lệch giữa các fold (0.011 và 0.008), nên hai mô hình coi như ngang nhau. Theo quy tắc chọn bằng $R^2$ validation,
\textbf{random forest} là mô hình cuối cùng.

\reportfig{fig09_hyperparameters.png}{$R^2$ train và validation theo \texttt{max\_depth} của cây và $k$ của kNN. Cây càng
sâu thì $R^2$ train càng tăng (0.905 khi không giới hạn), nhưng $R^2$ validation cao nhất là 0.789 ở \texttt{max\_depth}
= 10 và giảm còn 0.774 khi không giới hạn: cây quá sâu bị overfitting. kNN đạt validation cao nhất 0.768 ở $k = 10$, vẫn
kém random forest.}{fig:tab-hyper}

\subsection{One-hot hay target encoding cho dòng xe?}\label{sec:tab-enc}

Để kiểm tra lựa chọn mã hóa ở phần EDA, nhóm giữ nguyên mọi thứ khác và chỉ đổi cách mã hóa bốn cột nhiều mức: chỉ
one-hot (mỗi mức từ 5 tin một cột), chỉ target encoding, hoặc cả hai (cách đang dùng).

\begin{table}[H]\centering
\caption{$R^2$ validation và sai số \% median theo cách mã hóa bốn cột nhiều mức.}\label{tab:tab-enc}
\begin{tabular}{lrrr}
\toprule
Cách mã hóa & Số cột & Hồi quy tuyến tính & Random forest \\
\midrule
One-hot & 368 & 0.750 (20.1\%) & 0.848 (12.3\%) \\
Target encoding & 50 & 0.744 (23.0\%) & 0.856 (12.5\%) \\
Cả hai & 167 & 0.749 (21.3\%) & \textbf{0.858 (12.1\%)} \\
\bottomrule
\end{tabular}
\end{table}

\reportfig{fig10_encoding.png}{$R^2$ validation theo cách mã hóa. Với random forest, target encoding (0.856) tốt hơn
one-hot (0.848) mà chỉ cần 50 cột thay vì 368; dùng cả hai cao nhất (0.858). Với hồi quy tuyến tính, ba cách gần như
bằng nhau (0.744--0.750).}{fig:tab-enc}

One-hot phải gộp các mức ít tin (dòng xe + dung tích hiếm, quận ít tin) vào một cột ``hiếm'', còn target encoding giữ
được mức giá riêng của từng mức. Với hồi quy tuyến tính, one-hot nhỉnh nhất vì cho mỗi dòng xe một hệ số riêng. Nhóm giữ
cách ``cả hai'' cho mọi mô hình vì cách này tốt nhất với mô hình mạnh nhất và chỉ kém one-hot 0.001 với hồi quy tuyến
tính.

\reportfig{fig11_model_comparison.png}{$R^2$ validation của các mô hình (cột) và $R^2$ cross-validation trên train (vạch
đen, trung bình ± độ lệch giữa các fold). Random forest (0.858) và gradient boosting (0.856) chênh nhau ít hơn độ lệch
giữa các fold. So với baseline median theo dòng xe, sai số \% median giảm từ 30.0\% xuống 12.1\%.}{fig:tab-models}

\section{Đánh giá trên tập test}

\begin{table}[H]\centering
\caption{Kết quả trên tập test (2 374 tin), $R^2$ tính trên $\log_{10}$(giá). Hàng in đậm là mô hình cuối cùng, chọn
theo $R^2$ validation.}\label{tab:tab-test}
\begin{tabular}{lrrrr}
\toprule
Mô hình & $R^2$ & Sai số \% median & Lệch $\le$ 10\% & MAE (triệu) \\
\midrule
Baseline: giá trung bình & 0.000 & 49.6\% & 9.8\% & 15.75 \\
Baseline: median theo dòng xe & 0.574 & 28.3\% & 19.1\% & 9.57 \\
Hồi quy tuyến tính & 0.738 & 20.9\% & 26.4\% & 7.44 \\
kNN ($k=10$) & 0.781 & 19.3\% & 28.2\% & 6.88 \\
Cây quyết định (sâu 12) & 0.778 & 15.0\% & 36.4\% & 6.04 \\
\textbf{Random forest} & \textbf{0.851} & \textbf{11.6\%} & \textbf{45.4\%} & \textbf{4.72} \\
Gradient boosting & 0.853 & 12.7\% & 42.2\% & 4.98 \\
\bottomrule
\end{tabular}
\end{table}

Random forest đạt $R^2$ 0.851 trên test, gần với validation (0.858) và cross-validation (0.859). Một nửa số xe được đoán
lệch không quá 11.6\% giá rao, so với 28.3\% nếu chỉ tra giá median của dòng xe, tức sai số còn khoảng 40\%. Gradient
boosting có $R^2$ test nhỉnh hơn (0.853) nhưng sai số \% median cao hơn (12.7\%), nhất quán với kết luận ``ngang nhau'' ở
validation.

\reportfig{fig12_pred_vs_true.png}{Giá dự đoán so với giá rao thật trên test (thang log). 68.7\% số tin được đoán lệch
không quá 20\% (dải xám), ở mọi mức giá. Chỉ 217 tin (9.1\%) lệch quá 50\%; 71\% trong số đó là đoán cao hơn giá rao và
45\% thuộc 20\% tin rẻ nhất.}{fig:tab-pred}

\reportfig{fig13_feature_importance.png}{15 feature quan trọng nhất của random forest (các cột one-hot của cùng một
feature được cộng lại). Thông tin ``đây là xe gì'' (dòng xe, hãng, dung tích, loại xe) chiếm 70.7\% độ quan trọng, đời
xe và số km 14.3\%, nơi bán và người bán 9.9\%, khớp với phần EDA.}{fig:tab-imp}

\subsection{Mô hình sai ở đâu?}

\begin{table}[H]\centering
\caption{Sai số của random forest trên test theo mức giá và theo nơi bán, người bán. Cột cuối là trung vị của tỉ số giá
dự đoán / giá thật; lớn hơn 1 nghĩa là mô hình thường đoán cao hơn giá rao.}\label{tab:tab-err}
\begin{tabular}{lrrr}
\toprule
Nhóm & Số tin & Sai số \% median & Dự đoán / thật \\
\midrule
Giá 2--9 triệu & 483 & 19.7\% & 1.16 \\
Giá 9--16 triệu & 478 & 14.6\% & 1.05 \\
Giá 16--24 triệu & 463 & 10.2\% & 0.97 \\
Giá 24--38 triệu & 475 & 9.0\% & 0.96 \\
Giá 38--346 triệu & 475 & 7.6\% & 0.96 \\
\midrule
Hà Nội, cá nhân & 103 & 17.2\% & 0.97 \\
Hà Nội, cửa hàng & 346 & 12.4\% & 1.02 \\
TP.HCM, cá nhân & 527 & 15.6\% & 1.00 \\
TP.HCM, cửa hàng & 1 398 & 9.9\% & 0.99 \\
\bottomrule
\end{tabular}
\end{table}

Sai số giảm dần theo mức giá, từ 19.7\% ở 20\% tin rẻ nhất xuống 7.6\% ở 20\% tin đắt nhất. Ở nhóm rẻ nhất, mô hình
thường đoán cao hơn giá rao (tỉ số 1.16): với xe vài triệu, chênh một hai triệu đã là nhiều phần trăm. Năm tin bị đoán sai
nhiều nhất đều có giá rao 2--4.3 triệu nhưng được đoán 8--21 triệu, ví dụ một chiếc Honda Wave 2019 đi 332 km rao 2.5
triệu. Đây nhiều khả năng là tiền cọc hoặc giá mồi mà luật bỏ giá thấp vô lý ở bước tiền xử lý không bắt được. Theo người bán, nhóm khó đoán nhất là cá nhân ở Hà Nội (17.2\%), nhưng tỉ số dự đoán / giá thật của cả bốn nhóm đều
trong khoảng 0.97--1.02, nên mô hình không đoán lệch hẳn về một phía ở nhóm nào.

\section{Dự đoán giá xe mới}

Hàm \texttt{predict\_price()} nhận thông tin như khi đăng một tin trên Chợ Tốt, đi qua đúng các bước tạo feature và
pipeline của mô hình cuối cùng; cột nào không khai thì pipeline tự điền. Xe thuộc dòng ngoài phạm vi bị từ chối thay vì
trả về một con số không đáng tin. Bảng~\ref{tab:tab-cases} thay đổi lần lượt từng thông tin của cùng một chiếc Honda
Vision.

\begin{table}[H]\centering
\caption{Giá dự đoán khi thay đổi từng thông tin của một chiếc Honda Vision 2021, 20 000 km, người bán cá nhân ở Quận Gò
Vấp.}\label{tab:tab-cases}
\begin{tabular}{lr}
\toprule
Thay đổi & Giá dự đoán (triệu) \\
\midrule
Xe gốc & 24.5 \\
Đời 2017 thay vì 2021 & 17.6 \\
Đi 60 000 km thay vì 20 000 km & 21.7 \\
Bán ở Hà Nội (Cầu Giấy), cá nhân & 22.2 \\
Bán ở Hà Nội (Nam Từ Liêm), cửa hàng & 14.4 \\
Honda SH 150i thay vì Vision & 73.8 \\
\bottomrule
\end{tabular}
\end{table}

Đổi sang đời 2017 làm giá giảm 28\%, khoảng 8\% mỗi năm, gần với mức 7\% mỗi năm của Vision ở
Hình~\ref{fig:tab-dep}. Bán qua cửa hàng ở Nam Từ Liêm cho giá thấp hẳn (khoảng 0.6 lần giá ở TP.HCM): mô hình đã học
lại đúng hiện tượng ở Hình~\ref{fig:tab-city}. Đổi dòng xe sang SH thì giá tăng gấp ba, đúng như dòng xe là yếu tố quyết
định.

\section{Kết luận}

Trên 65 dòng xe phổ biến (82.8\% số tin sau khi làm sạch), random forest giải thích 85\% biến động của $\log$(giá) trên
tập test và đoán một nửa số xe lệch không quá 11.6\% giá rao, so với 28.3\% của baseline tra giá median theo dòng xe.
Gradient boosting đạt kết quả tương đương. Ba câu hỏi đặt ra được trả lời như sau:

\begin{itemize}
  \item \textbf{Lỗi của dữ liệu tự thu thập:} 74 tin giá rác hoặc tin cần mua, 517 tin đăng lại, 10.8\% tin xe đã dùng
        khai số km ảo, 3 cột hằng và tọa độ chỉ ở mức phường.
  \item \textbf{Mất giá:} khoảng 5--9\% mỗi năm đời xe với 4 dòng phổ biến nhất; xe đi trên 80 000 km còn 61--88\% giá
        của xe đi dưới 5 000 km.
  \item \textbf{Hà Nội so với TP.HCM:} cá nhân ở Hà Nội rao bằng 0.90 lần giá TP.HCM, cửa hàng chỉ 0.62 lần. Chênh lệch
        chủ yếu đến từ một nhóm cửa hàng, không phải toàn thị trường.
\end{itemize}

\subsection*{Hạn chế}
\begin{itemize}
  \item \textbf{Giá rao không phải giá bán.} Nhóm cửa hàng rao rất thấp ở Hà Nội và các tin rẻ bất thường ở test cho
        thấy một phần giá rao có thể không phải giá giao dịch thật.
  \item \textbf{Điểm test có thể hơi lạc quan.} Trước khi viết bản notebook cuối, nhóm đã thử vài phạm vi dữ liệu và vài
        cách mã hóa, và có xem điểm test của các lần thử đó; quyết định giới hạn vào 65 dòng xe phổ biến được đưa ra sau
        các lần thử này. Mức lạc quan có lẽ nhỏ, vì $R^2$ validation (0.858) và cross-validation (0.859) gần với test
        (0.851).
  \item \textbf{Luật bỏ giá thấp vô lý dùng nhãn:} 5 tin ở validation và test bị bỏ dựa trên chính giá rao của chúng.
  \item \textbf{Thiếu thông tin tình trạng xe.} Trầy xước, máy móc, giấy tờ nằm trong ảnh và phần mô tả mà nhóm không thu
        thập để bảo vệ thông tin cá nhân. Số km do người bán tự khai, 11\% số tin bị coi là thiếu số km.
  \item \textbf{Một lần thu thập và chỉ 65 dòng xe.} Xe bán nhanh có thể đã bị gỡ trước ngày thu thập; các dòng hiếm
        (11.6\% số tin) và \texttt{Dòng khác} (5.6\%) nằm ngoài phạm vi.
\end{itemize}

\subsection*{Hướng mở rộng}
Đưa phần mô tả tin vào mô hình bằng TF-IDF hoặc embedding tiếng Việt (như bài Text) và ảnh xe bằng mạng pretrained (như
bài Image) để có thông tin về tình trạng xe; thu thập nhiều đợt để theo dõi giá theo thời gian và biết tin nào được gỡ
sớm; dự đoán một khoảng giá thay vì một con số; tinh chỉnh tham số bằng \texttt{GridSearchCV} và mở rộng sang các tỉnh
thành khác.

\begin{thebibliography}{9}
\bibitem{tab-chotot} Chợ Tốt. \textit{Mua bán xe máy cũ}. \url{https://xe.chotot.com/mua-ban-xe-may}, truy cập ngày
  16/11/2026.
\bibitem{tab-sklearn} F. Pedregosa và cộng sự. \textit{Scikit-learn: Machine Learning in Python}. Journal of Machine
  Learning Research, 12, 2011, tr. 2825--2830.
\bibitem{tab-te} D. Micci-Barreca. \textit{A Preprocessing Scheme for High-Cardinality Categorical Attributes in
  Classification and Prediction Problems}. ACM SIGKDD Explorations Newsletter, 3(1), 2001, tr. 27--32.
\bibitem{tab-rf} L. Breiman. \textit{Random Forests}. Machine Learning, 45(1), 2001, tr. 5--32.
\end{thebibliography}
